% 付録C — 検証と制限
\chapter{検証の読み方とモデルの制限}\label{app:validation}

テストやデモが示す性質の種類を区別してください。決定論的な検査は特定入力に対する恒等式や真理値表を確かめます。統計検査は有限ショットの推定値であり、乱数シード、ショット数、判定閾値とともに解釈します。距離計算が \texttt{shortest\_graphlike\_error()} を使う場合、その値は graphlike 誤りモデルの範囲での上界です。未分解の全誤り機構について別途条件を確認できない限り、符号の厳密距離や一般の FT 性の証明とは呼べません。

公開版には各 Stage のテスト一式と実行例を含めています。テストの存在だけで全項目の実行成功を意味するわけではありません。再現する場合は付録Aのコマンドを実行し、実際の終了コードと出力を記録してください。本文中で個別に示す数値例も、特定の設定・実装に対する結果として扱ってください。

Stage 01 は安定化子を MPP で直接測る現象論モデルであり、局所 ancilla 回路や hook 誤りを表しません。Stage 02 は回路レベルの誤りを導入しますが、任意距離の一般化や実機の全ノイズ過程を含みません。Stage 03 は $d=3$ と 9 個のデータ量子ビットに固定された統合シミュレータです。残差モデルは Stim から得る観測量フリップを中心に構成され、$X_L$ と $Z_L$ の相関を一般に完全に表すものではありません。

したがって、この教材の実行結果から大規模 FTQC の閾値、実機性能、一般の論理誤り率を直接推定することはできません。数値を比較する場合は、回路、ノイズ仮定、デコーダ、距離、サンプリング条件を併記してください。
